%% hopfion-framework/foundations/pentagon.tex
%% Sources: main_paper_foundational.tex, main_paper1.tex
%% The Pentagon Theorem: k+2=5 and all its consequences
%% Auto-generated from project files. Edit source papers to update.
%%
%% Status tags: \status{published}, \status{open}, \status{conjectured}
%% \source{PaperN, label} — where the proof lives
%% \residual{X\%} or \residual{exact} — numerical accuracy


%════════════════════════════════════════════════════════════════════
% Pentagon Theorem — Foundational Companion
%════════════════════════════════════════════════════════════════════
% The four-part theorem from which everything derives
\begin{theorem}[The Pentagon Theorem]
\label{F:thm:main}
In the density-feedback Faddeev--Niemi Hopfion condensate:

\begin{enumerate}[label=\emph{(\roman*)}]
  \item \emph{(Uniqueness of $k=3$.)}
    The WZW level is $k=3$, excluded at all other positive integer
    values by QED at $\sim1.5\times10^7\,\sigma$ (Proposition~\ref{F:prop:route1}).

  \item \emph{(Three generations from primary truncation.)}
    $\mathrm{SU}(2)_3$ has exactly three non-vacuum primaries
    $j = \tfrac{1}{2}, 1, \tfrac{3}{2}$, corresponding to the electron,
    muon, and tau.
    A fourth primary $j=2$ does not exist (Proposition~\ref{F:prop:route3}).

  \item \emph{(Golden ratio as quantum dimension.)}
    Every non-vacuum primary has quantum dimension $\vp$ or $1$:
    $d_{1/2} = d_1 = \vp$, $d_{3/2} = 1$ (Proposition~\ref{F:prop:route2}
    and Remark~\ref{F:rem:both_phi}).

  \item \emph{(Common origin: $k+2=5$.)}
    All three results are consequences of the single equation $k+2=5$.
    The connection is:
    \begin{itemize}
      \item $k+2=5$ $\Rightarrow$ $d_{1/2} = 2\cos(\pi/5) = \vp$
        (pentagon half-angle);
      \item $k+2=5$ $\Rightarrow$ $j_{\max}=k/2=3/2$
        $\Rightarrow$ three non-vacuum primaries;
      \item $k+2=5$ $\Rightarrow$ $k=3$ is the level selected by QED.
    \end{itemize}
\end{enumerate}
\end{theorem}
\status{published}
\source{F:thm:main}

\begin{proposition}[Route~1: QED exclusion]
\label{F:prop:route1}
At $k=3$, equation~\eqref{P3:eq:alpha_formula} gives
$\alpha^{-1} = 137.030$, an error of $0.0042\%$ from the experimental
value $\alpha^{-1}_{\mathrm{QED}} = 137.035999084$.
At $k=4$, $\alpha^{-1} = 136.871$, an error of $0.120\%$.
The experimental precision of $\alpha$ is $\delta\alpha/\alpha
\approx 6\times10^{-11}$~\cite{Morel2020}, giving the exclusion
\begin{equation}
  |k=4 \text{ error}| \;\approx\; 1.5\times10^7\,\sigma.
  \label{F:eq:sigma_exclusion}
\end{equation}
\end{proposition}
\status{published}
\source{F:prop:route1}

\begin{proposition}[Route~2: Quantum dimension uniqueness]
\label{F:prop:route2}
The quantum dimension of the muon primary $j=1/2$ equals the golden ratio,
\begin{equation}
  d_{1/2}\big|_{k=3} = \vp,
  \label{F:eq:qdim_phi}
\end{equation}
if and only if $k=3$.
At no other positive integer level does $d_{1/2} = \vp$.
\end{proposition}
\status{published}
\source{F:prop:route2}

\begin{remark}[Both $j=1/2$ and $j=1$ share quantum dimension $\vp$]
\label{F:rem:both_phi}
At $k=3$, \emph{both} $j=1/2$ and $j=1$ have quantum dimension $\vp$:
\begin{equation}
  d_{1/2}\big|_{k=3}
  = d_{1}\big|_{k=3}
  = \vp.
  \label{F:eq:both_phi}
\end{equation}
This follows since $\sin(3\pi/5) = \sin(\pi - 2\pi/5) = \sin(2\pi/5)$,
so $d_1 = \sin(3\pi/5)/\sin(\pi/5) = \sin(2\pi/5)/\sin(\pi/5) = \vp$.
The tau primary $j=3/2$ has $d_{3/2} = \sin(4\pi/5)/\sin(\pi/5)
= \sin(\pi/5)/\sin(\pi/5) = 1$: the tau is its own antiparticle
under fusion (it has unit quantum dimension).
This is the quantum group analogue of the Majorana condition,
and it holds at $k=3$ only.
\end{remark}
\status{published}
\source{F:rem:both_phi}

\begin{proposition}[Route~3: No fourth generation]
\label{F:prop:route3}
A fourth lepton generation would require primary $j_4=2$.
This primary does not exist in $\mathrm{SU}(2)_3$:
\begin{equation}
  j_4 = 2 \;>\; \frac{k}{2} = \frac{3}{2}.
  \label{F:eq:j4_excluded}
\end{equation}
This is a hard algebraic exclusion, independent of any numerical input.
\end{proposition}
\status{published}
\source{F:prop:route3}

\begin{remark}[Strength of the three routes]
\label{F:rem:route_strength}
The three routes are logically independent: each would individually
exclude $k\neq3$. But they differ in character:
Route~1 is quantitative (requires the numerical value of $\alpha$);
Route~2 is algebraic-exact (requires only that $\vp$ be the golden ratio);
Route~3 is a hard algebraic non-existence statement (requires only the
Kac--Moody unitarity condition, which is a theorem in mathematics,
not an empirical input).
Route~3 is therefore the strongest: even if the fine structure constant
formula~\eqref{F:eq:alpha_formula} were corrected by higher-order terms,
and even if the quantum dimension identity~\eqref{F:eq:qdim_phi} were
modified by quantum corrections, Route~3 would still apply.
\end{remark}
\status{published}
\source{F:rem:route_strength}

\begin{proposition}[$Q=10$ and $k=3$ as McKay invariants]
\label{F:prop:mckay}
Both $Q=10$ and $k=3$ are invariants of $\twoI/\Ehat$:
\begin{itemize}
  \item $k=3$: the WZW level is determined by the embedding
    $\twoI \hookrightarrow \mathrm{SU}(2)$ and the level-3 truncation
    that selects the icosahedrally symmetric primary sector.
  \item $Q=10$: the Hopf charge equals the number of $\twoI$-orbits
    on the condensate configuration space, fixed by the icosahedral
    symmetry of the FN functional~\cite{Paper1}.
\end{itemize}
The precise algebraic relationship is $Q = 2(k+2)$
(Paper~I~\cite{Paper1}, Theorem~5.9 of that paper: $\mathrm{ord}(q_{2I})=2(k+2)=10$),
and $k\,h(E_8) = Q\times 9 = 90$
(Paper~V~\cite{Paper5}, Theorem~11.8 of that paper: $3\times30=90=Q\times9$).
\end{proposition}
\status{published}
\source{F:prop:mckay}

\begin{corollary}[The golden ratio is not coincidental]
\label{F:cor:phi_origin}
The golden ratio $\vp$ appears throughout the condensate framework
--- in $\lam=\vp^6$, $c_s=1/\vp$, $m_n=m_0\vp^{-2n}$,
$\det(\mathcal{L}_*)=\vp^{-1/6}$, $d_{1/2}=\vp$ ---
because it is the half-angle cosine of the pentagon,
and the condensate's WZW sector is governed by $k+2=5$.
It is not that $\vp$ was found empirically and then built into the theory;
it is that $k+2=5$ forces $\vp$ everywhere $\mathrm{SU}(2)_3$
quantum dimensions appear.
\end{corollary}
\status{published}
\source{F:cor:phi_origin}

\begin{remark}[Three descriptions of the same constraint]
\label{F:rem:three_descriptions}
The following are not three independent reasons why there are three
\emph{lepton} generations: they are three descriptions of the same constraint.
\begin{enumerate}[label=(\arabic*)]
  \item Three non-vacuum primaries of $\mathrm{SU}(2)_3$.
  \item Three non-real Hurwitz division algebras $\CC$, $\HH$, $\OO$.
  \item Sedenions (the $k=4$ Cayley--Dickson algebra) lose the
    division property: the Hurwitz tower terminates at $\OO$.
\end{enumerate}
The algebraic reason for all three is the same: the sequence
$\RR\to\CC\to\HH\to\OO$ loses one algebraic property at each step
(ordering, commutativity, associativity), and the next step would lose
alternativity, destroying the norm identity.
This is encoded in $k+2=5$ through the Verlinde~\cite{Verlinde1988} formula:
the fusion ring of $\mathrm{SU}(2)_3$ is the representation ring of
the quantum group $\mathcal{U}_q(\mathfrak{su}(2))$ at $q=e^{2\pi i/5}$,
a fifth root of unity, whose truncation at $j=3/2$ mirrors
the Hurwitz termination.

For \emph{quarks}, the generation count arises from a different mechanism
within the same $\twoI/\Ehat$ structure: the six SM quarks occupy six of the
eight $E_8$ Coxeter exponents, with the two unoccupied exponents $\{23,29\}$
predicting no fourth quark generation (Paper~V~\cite{Paper5},
Proposition~8 therein).
This is $E_8$ Coxeter saturation, not $\mathrm{SU}(2)_3$ primary truncation,
as established in Paper~VI~\cite{Paper6}, Remark~5.6 therein.
\end{remark}
\status{published}
\source{F:rem:three_descriptions}


%════════════════════════════════════════════════════════════════════
% Key Equations — Foundational
%════════════════════════════════════════════════════════════════════
% Pentagon identity, quantum dimensions, generation count
\begin{equation}
  k + 2 = 5,
  \label{F:eq:k_plus_2}
\end{equation}

\begin{equation}
  j_g = \frac{g}{2}, \qquad g = 1, 2, 3.
  \label{F:eq:generation_primaries}
\end{equation}

\begin{equation}
  d_j \;=\; \frac{\sin\!\bigl((2j+1)\pi/(k+2)\bigr)}
                  {\sin\!\bigl(\pi/(k+2)\bigr)}.
  \label{F:eq:qdim}
\end{equation}

\begin{equation}
  h_j \;=\; \frac{j(j+1)}{k+2}.
  \label{F:eq:conformal_weight}
\end{equation}


%════════════════════════════════════════════════════════════════════
% Pentagon results — Paper I
%════════════════════════════════════════════════════════════════════
% k=3 uniqueness, icosahedral symmetry, Q=10
\begin{equation}
  D_0 = (r-R_0)^2 + z^2, \qquad
  A   = \frac{1}{D_0} + \frac{1}{r^2}, \qquad
  \mathrm{kern} = |\nabla f|^2 + \sin^2\!f\cdot A.
  \label{P1:eq:kern}
\end{equation}

